\originalchapter{流形与长度几何}\label{ch:manifolds}
\emph{坐标图、光滑映射与逆函数论证依据Tomasz S. Mrowka的18.965 Fall 2004第1、2、4讲; 路径长度及内在距离参照Paul Seidel的18.950 Fall 2008第4章Lecture 36、39. 有限维局部构造、一般已给定黎曼度量后的距离拓扑证明及配套题解为编者补充. 本章不涵盖这些课程的全部流形、曲率或测地线理论.}

\originalsection{局部欧氏空间}
第8章将流形作为分析的发展方向. 现在可以将局部坐标、全局拓扑和距离的关系具体建立起来. 一个圆周可以用小开弧的角度描述, 但全局没有一张实区间坐标图; 坐标重叠的变换使多张局部描述能够相容.
\begin{definition}
一个$n$维拓扑流形是Hausdorff、第二可数且每点有开邻域同胚于$\R^n$开子集的空间. 坐标图为$(U,\varphi)$, 其中$U$开且$\varphi:U\to\varphi(U)\subseteq\R^n$为同胚, $\varphi(U)$须开. 若图族覆盖空间, 且转移映射$\psi\circ\varphi^{-1}$在重叠坐标域上均光滑, 称为光滑图册. 指定相容光滑图册的流形称为光滑流形.
\end{definition}
维数$n$在本册的定义中固定; 本册不证明维数不变定理. 有边界流形把坐标像改为半空间$\{x_n\ge0\}$的相对开集, 边界点映到$x_n=0$. 以下光滑微分讨论默认无边界, 带边界情形需额外约定延拓的光滑性.

\begin{proposition}\label{thm:manifold-metrizable}
拓扑流形局部紧、局部道路连通且可度量化. 连通流形道路连通.
\end{proposition}
\begin{proof}
在坐标像内以点为中心取一个闭球, 半径足够小使其留在坐标域. 逆像是紧邻域, 因闭球紧而图逆映射连续, 故局部紧. 更小的开球逆像道路连通, 并可包含于任意给定邻域, 故局部道路连通. 局部紧Hausdorff空间正则, 用引理\ref{lem:local-compact-shrink}; 第二可数正则空间可度量化, 用定理\ref{thm:metrization}. 道路分支在局部道路连通空间中开, 若连通只能有一个.
\end{proof}
Hausdorff和第二可数不能只因有局部坐标就省去. 例如两份实直线将全部非零点按相同坐标识别, 却保留两个零点, 形成“双原点直线”. 每点有实区间坐标, 但两个原点的任意邻域均相交, 不是Hausdorff, 因而不符合本书流形定义.

\begin{proposition}
每个光滑图册唯一包含于一个极大相容图册. 构造不需要另用Zorn引理.
\end{proposition}
\begin{proof}
把所有与原图册每张图都相容的图加入. 任取两张新图, 在它们重叠的每一点, 选一张原图覆盖该点; 两新图之间的转移映射在该点附近为通过原图的两个光滑转移映射的复合, 故光滑, 反向也一样. 因而新族仍为图册. 任何与原图册相容的图已经在新族内, 故极大且唯一.
\end{proof}

\originalsection{球面与射影空间的坐标}
\begin{example}
写出$S^n$与$\mathbb RP^n$的光滑图册, 并核对转移映射.
\end{example}
\begin{solution}
在$S^n\subseteq\R^n\times\R$中记点$(u,t)$. 去掉北极时取$\varphi_N(u,t)=u/(1-t)$, 去掉南极时取$\varphi_S(u,t)=u/(1+t)$. 逆映射分别为
\[
 x\longmapsto\left(\frac{2x}{1+|x|^2},\frac{|x|^2-1}{1+|x|^2}\right),\qquad
 x\longmapsto\left(\frac{2x}{1+|x|^2},\frac{1-|x|^2}{1+|x|^2}\right).
\]
重叠处$x\ne0$, 转移为$x\mapsto x/|x|^2$, 光滑且自逆. 球面为欧氏子空间, Hausdorff且第二可数, 两图覆盖, 所以确为光滑流形.

射影空间沿用第10章的球面对径商拓扑, 将一个对径类记为$[x_0:\cdots:x_n]$, 即非零向量所在的直线. 每条直线恰有两个单位向量, 所以这只是同一商空间的另一种记法. 开集$U_i=\{[x]:x_i\ne0\}$上取$\varphi_i([x])=(x_j/x_i)_{j\ne i}$. 比值在对径点上相同, 由商泛性质连续; 逆映射插入$x_i=1$, 再归一化并取对径类, 连续. 图间转移通过非零坐标相除, 因而在其开放定义域光滑. $\mathbb RP^n$已证紧Hausdorff, 有限张第二可数图的基之并为可数全局基, 故满足流形条件. 写比值时必须删去第$i$项, 否则坐标多出恒等于$1$的一项.
\end{solution}

\begin{definition}
光滑流形间映射$f:M\to N$称光滑, 若连续且所有局部坐标表达$\psi\circ f\circ\varphi^{-1}$光滑. 双射$f$及其逆均光滑时称为微分同胚.
\end{definition}
只需对覆盖定义域、目标域的一组图检查. 换图后表达式在两侧复合光滑转移映射, 由链式法则仍光滑. 同胚只约束连续结构, 微分同胚还约束指定光滑结构; 在$\R$上$x\mapsto x^3$为同胚, 但其逆$x^{1/3}$在零点不可微, 不是通常光滑结构下的微分同胚.

\originalsection{逆函数与正则水平集}
本节用多元微积分的导数、链式法则及矩阵运算作为先修. 关键存在性由第6章收缩映射原理证明, 使局部坐标的构造不只停在定理名称.
\begin{theorem}[有限维逆函数定理]\label{thm:inverse-function}
设$F:O\subseteq\R^n\to\R^n$为$C^1$映射, $O$开, $DF(a)$可逆. 则有$a$的开邻域$U$及$F(a)$的开邻域$V$, 使$F:U\to V$为$C^1$微分同胚, 且$D(F^{-1})(y)=DF(F^{-1}(y))^{-1}$. 若$F$光滑, 逆也光滑.
\end{theorem}
\begin{proof}
平移并复合可逆线性映射, 可令$a=0$, $F(0)=0$, $DF(0)=I$. 写$F(x)=x+G(x)$. 由导数连续性选$r>0$, 使闭球$\overline B_r\subseteq O$且其上$\lVert DG\rVert\le1/2$. 沿线段积分导数得$|G(x)-G(x')|\le|x-x'|/2$. 对$|y|<r/2$, 映射$T_y(x)=y-G(x)$将闭球映入自身, 因为$G(0)=0$且$|T_y(x)|\le|y|+r/2<r$. 它为收缩, 故有唯一固定点$\phi(y)$满足$F(\phi(y))=y$.

两个固定点方程相减, 得$|\phi(y)-\phi(y')|\le2|y-y'|$, 所以$\phi$连续. 还由$|\phi(y)|\le2|y|<r$得其值在开球内. 令$V=B_{r/2}$, $U=B_r\cap F^{-1}(V)$, 则$U$开, 每个$y\in V$恰对应一个$x\in U$, 所以$F:U\to V$与$\phi$互逆.

在闭球内$DF=I+DG$均可逆: 级数$\sum_{k\ge0}(-DG)^k$在算子范数下收敛, 与$I+DG$相乘的有限部分和余项趋零, 得逆矩阵且逆范数不超过$2$. 对$x=\phi(y)$, $h=\phi(y+k)-\phi(y)$, 有$k=DF(x)h+o(|h|)$, 而$|h|\le2|k|$. 乘逆矩阵得到$h=DF(x)^{-1}k+o(|k|)$, 故$D\phi(y)=DF(\phi(y))^{-1}$. 右侧连续, 因而逆为$C^1$. 若$F$光滑, 矩阵求逆在可逆矩阵集合中光滑, 可用伴随矩阵除以非零行列式核对; 由此公式和链式法则逐阶归纳, 得$\phi$光滑.
\end{proof}

\begin{theorem}[正则水平集]\label{thm:regular-level}
设$f:O\subseteq\R^n\to\R^m$光滑, $O$开, $m\le n$, 且对每个$x\in f^{-1}(c)$, $Df(x)$满射. 则非空水平集$f^{-1}(c)$按子空间拓扑为光滑$(n-m)$维流形.
\end{theorem}
\begin{proof}
在某点$a$, 满秩矩阵有一个$m\times m$可逆子块. 重排坐标为$(u,v)\in\R^{n-m}\times\R^m$, 使$D_vf(a)$可逆. 映射$\Phi(u,v)=(u,f(u,v))$的导数为块三角矩阵, 对角块为$I$与$D_vf(a)$, 故可逆. 由逆函数定理, $\Phi$在$a$附近为微分同胚. 水平集在此坐标下恰为第二组坐标等于$c$的切片, 第一组坐标$u$形成开集, 给局部图. 两图间转移是环境光滑变换在该切片上的限制, 因而光滑. 水平集作为欧氏子空间Hausdorff且第二可数, 具备全部流形条件.
\end{proof}
例如$f(x)=|x|^2$, $c=1$时$Df(x)h=2x\cdot h$满射, 得球面为光滑流形. 零水平集未满足正则条件; 它仍可能碰巧是其他维数的流形, 所以满秩是充分条件, 不能称为所有水平集成为流形的必要条件. 曲线$xy=0$在原点呈十字, 则确实不局部同胚于实区间: 删去原点的小邻域有四个分支, 区间删去一点仅两个.

\originalsection{切空间与微分}
\begin{definition}
在光滑流形点$p$, 将过$p$的光滑曲线按在一张局部图中有相同的零时刻导数分类, 得切向量. 这些类组成切空间$T_pM$. 光滑映射$f:M\to N$的微分$df_p:T_pM\to T_{f(p)}N$将曲线类$[\gamma]$送到$[f\circ\gamma]$.
\end{definition}
这里曲线定义在含$0$的小开区间, $\gamma(0)=p$. 换坐标时速度由可逆矩阵$D(\psi\circ\varphi^{-1})(\varphi(p))$变换, 所以曲线是否同类与图无关. 每个坐标速度$v\in\R^n$由曲线$\varphi^{-1}(\varphi(p)+tv)$实现, 因而一张图给出$T_pM\cong\R^n$. 加法和数乘在图中定义, 可逆线性坐标变化使之独立于图. 同理链式法则使$df_p$良定义且线性, 并给$d(gf)_p=dg_{f(p)}\circ df_p$.

\begin{proposition}\label{thm:tangent-level}
正则水平集$M=f^{-1}(c)\subseteq\R^n$的切空间为$T_pM=\ker Df(p)$, 其中环境速度用于识别切向量.
\end{proposition}
\begin{proof}
对留在$M$中的曲线,$f(\gamma(t))=c$, 求导得$Df(p)\gamma'(0)=0$. 反向用正则水平集证明中的坐标$\Phi=(u,f)$, 微分将$\ker Df(p)$映成第一坐标空间$\R^{n-m}\times\{0\}$. 在水平切片中取给定速度的直线, 经$\Phi^{-1}$得到$M$中的曲线, 环境速度正是原向量. 因而两个包含方向均成立.
\end{proof}
球面上$T_pS^n=\{v:p\cdot v=0\}$. 切空间是附在一点的线性空间, 不等于曲面附近的一块; 将它画为过$p$的仿射平面时, 实际画的是$p+T_pM$.

\originalsection{长度与内在距离}
\begin{definition}
黎曼度量是在每个$T_pM$上指定正定内积$g_p$, 在局部图中其矩阵系数$g_{ij}(x)$随$x$光滑. 分段$C^1$曲线$\gamma:[a,b]\to M$的长度为
$L_g(\gamma)=\int_a^b\sqrt{g_{\gamma(t)}(\gamma'(t),\gamma'(t))}\,\mathrm dt$,
其中对$C^1$曲线, $\gamma^{\prime}(t)$指坐标导数经$T_{\gamma(t)}M\cong\R^n$对应得到的切向量. 在分段点分别积分相加, 有限个端点的速度取值不影响积分.
\end{definition}
坐标变化使切向量和度量矩阵按相反方向变换, 内积值不变, 所以长度与图无关. 保端点的单调满射$C^1$重参数在可逆区间上由积分代换保持长度; 允许分段重参数或等待段时分段核对即可. 对嵌入欧氏空间的正则水平集, 环境内积限制在切空间给诱导度量. 参数化$F(u)$的坐标矩阵为$g_{ij}=\partial_iF\cdot\partial_jF$, 正定性来自参数导数的单射性.

\begin{theorem}\label{thm:riemannian-distance}
连通光滑流形给定黎曼度量后, $d_g(p,q)=\inf_\gamma L_g(\gamma)$是有限度量, 并生成原流形拓扑, 下确界取遍连接两点的分段$C^1$曲线. 本定理不声称下确界必被某条曲线取得.
\end{theorem}
\begin{proof}
坐标小凸球中的点可以用坐标直线连接, 对应光滑曲线. 按可用分段光滑曲线连通划分, 每类开, 连通性使只有一类, 所以下确界有限. 反向道路给对称性, 拼接并按任意小误差选近乎最短曲线给三角不等式, 常值曲线给$d_g(p,p)=0$.

须另外证明不同点距离正以及拓扑一致. 在$p$的一张坐标图中取中心为$p$的闭坐标球$\overline B_r$, 紧含于图域. 度量矩阵连续且正定, 在$\overline B_r\times S^{n-1}$上取最小、最大值, 得$0<c\le C<\infty$满足
$c|v|\le\sqrt{g_x(v,v)}\le C|v|$.
对留在此球内的曲线, 长度至少为$c$乘坐标端点距离. 对离开球的曲线, 到第一次碰到坐标球边界为止的长度至少$cr$; 首次碰边存在, 因曲线连续且闭球紧含于图域, 从中心离开开球必须先经过其边界. 因而对$q$在球内, 任意曲线长度至少$\min\{c|\varphi(q)-\varphi(p)|,cr\}$, 若$q$在球外则至少$cr$. 不同$q$的下界正. 球内直线又给$d_g(p,q)\le C|\varphi(q)-\varphi(p)|$.

上界说明坐标点趋于$p$时$d_g$趋零; 下界说明足够小$d_g$球包含于任意给定的小坐标球. 两种邻域系相互包含, 故拓扑一致. $n=0$时连通非空流形只有一点, 结论直接成立.
\end{proof}
这是内在距离: 路径须留在流形内. 在单位圆周上, 两点的内在距离为较短圆弧长度, 环境距离为弦长, 两者生成相同拓扑却数值不同. 任意图册给出的可度量化也不唯一指定黎曼度量; 本节定理从已经给定的度量出发, 不代替一般流形上黎曼度量存在性的单位分解证明.

\begin{example}
计算圆柱$M=\{(\cos\theta,\sin\theta,z)\}$的内在距离, 并解释局部坐标为何不能直接消除全局绕圈.
\end{example}
\begin{solution}
参数导数分别为$(-\sin\theta,\cos\theta,0)$与$(0,0,1)$, 故局部度量为$\mathrm d\theta^2+\mathrm dz^2$. 对两点取角代表$\theta_1,\theta_2$. 记覆盖$P(\theta,z)=(\cos\theta,\sin\theta,z)$, 任意连接曲线取提升初值$(\theta_1,z_1)$, 终点为$(\theta_2+2\pi k,z_2)$, $k\in\mathbb Z$. 在小开圆弧与高度区间上, $P$的逆支是光滑角度坐标. 紧参数区间有有限分割使每小段留在一张这样的图内, 并可细化原来的$C^1$分段; 提升在各段等于光滑逆支与原曲线复合, 故仍分段$C^1$. 若提升写为$(\theta(t),z(t))$, 原曲线速度为$(-\sin\theta\,\theta^{\prime},\cos\theta\,\theta^{\prime},z^{\prime})$, 范数平方为$(\theta^{\prime})^2+(z^{\prime})^2$, 因而两曲线长度相等. 欧氏平面曲线长度至少为端点距离, 等号可由直线实现, 投影回圆柱得
\[
 d_M(p,q)=\min_{k\in\mathbb Z}
 \sqrt{(\theta_2-\theta_1+2\pi k)^2+(z_2-z_1)^2}.
\]
整数最小值存在, 因表达式随$|k|\to\infty$趋无穷. 单一角坐标只能覆盖开圆弧, 整圈的整数差须由覆盖提升记录. 这是拓扑与长度计算共同作用的一个具体例子.
\end{solution}

\originalsection{练习与解答}
\begin{exercise}
证明$S^1\times S^1$为光滑二维流形, 写出标准乘积度量并计算两个角坐标点的距离.
\end{exercise}
\begin{solution}
两开圆弧图的积覆盖环面, 转移在各重叠连通分支上为角度加常数$2\pi k$, 因而光滑; Hausdorff和第二可数由有限积结论. 乘积度量为$\mathrm d\theta^2+\mathrm d\phi^2$. 由圆柱例题的逆支分段论证, 覆盖$P(\theta,\phi)=(e^{i\theta},e^{i\phi})$将连接曲线提升为$\R^2$中的分段$C^1$曲线, 且逐段速度平方为$(\theta^{\prime})^2+(\phi^{\prime})^2$, 长度保持. 固定提升初值后, 两终点差为$(\Delta\theta+2\pi m,\Delta\phi+2\pi n)$, 所以距离是这些欧氏范数的最小值, $m,n\in\mathbb Z$. 直线投影实现最小值. 这是平坦环面的内在几何, 不等于常见空间旋转环面从$\R^3$诱导的度量.
\end{solution}
\begin{exercise}
对$f(x,y,z)=x^2+y^2-z^2$, 证明每个非零水平集是光滑曲面; 分析零水平集在原点附近的拓扑.
\end{exercise}
\begin{solution}
$Df=(2x,2y,-2z)$仅在原点为零, 非零水平集不含原点, 故正则且二维. 零水平集为双圆锥. 删掉原点后分成$z>0$与$z<0$两部分, 两部分在任意原点邻域内都非空且不能在删点圆锥中连接. 若原点有二维圆盘坐标邻域, 缩小为坐标圆盘后删中心仍道路连通, 与圆锥的这两分支矛盾. 因而原点不是二维流形点; 其余点仍正则.
\end{solution}
\begin{exercise}
在开单位圆盘内采用欧氏诱导度量. 证明它的内在距离就是欧氏距离, 却不是完备空间. 说明这与局部欧氏性是否冲突.
\end{exercise}
\begin{solution}
圆盘凸, 线段连接给长度等于弦长, 任何曲线长度不小于弦长, 故两距离相同. 点列$(1-1/n,0)$($n\ge2$)为Cauchy列, 在平面极限为$(1,0)$但不属于开圆盘, 因而在圆盘内不收敛. 完备性是选定距离的全局性质, 局部坐标只描述邻域结构, 所以没有冲突. 本例也不包含关于完备黎曼流形上最短曲线存在的Hopf--Rinow定理.
\end{solution}
